\section{Fixed-model nonzero coordinates are exactly the original nonempty charts}
\label{sec:projective-native-coordinate-chart-support}

\begin{theorem}[Exact native support of every original generic coordinate]
\label{thm:projective-native-coordinate-chart-support}
\lean{EnvelopeProperty.ProjectiveNativeCoordinateChartSupport.coordinate_ne_zero_iff_genericPoint_mem, EnvelopeProperty.ProjectiveNativeCoordinateChartSupport.coordinate_ne_zero_iff_nonempty}
\leanok
\uses{thm:projective-chart-generic-transitions}
\source{bj22-trivval:sec:tc}
For an actual original morphism from integral $X$ to polynomial
projective space, fix one nonempty original coordinate chart $i$.
The literal original generic coefficient $c_{i,j}$ in native $K(X)$
is nonzero exactly when the original generic point maps to chart $j$,
and exactly when the original inverse image of chart $j$ is nonempty.
No nonempty-chart or support-matching premise for $j$ is supplied.
\end{theorem}
\begin{proof}
On the original $i$th chart, the actual function $X_j/X_i$ has principal
open precisely the overlap with the original $j$th chart. If its actual
pulled generic germ is nonzero, it is a unit in the actual native generic
stalk field. The native basic-open criterion therefore places the
original generic point in the principal open of this pulled section.
The original scheme pullback/basic-open identity and the actual
projective local-equation identity put its image in chart $j$.

Conversely membership of the original generic point makes its original
inverse chart nonempty, and the proved actual generic chart transition
forces its alternative pivot to be nonzero. A nonempty open of an
integral scheme contains its original generic point, while any open
containing that point is nonempty. This gives the second equivalence,
with all original coordinate functions and native scheme opens retained.
\end{proof}
